%!TEX root = ../main.tex
% =============================================================================
%  十、模型评价与推广            建议 2 页
%  评审在问：你诚实吗？能推广吗？
%  本节必须出现：① 创新五段式（缺陷—改动—证据—代价—边界，每条都要有消融/对照数字）
%                ② 局限四要素（指名假设编号 + 影响方向 + 所需数据 + 校准动作）
%                ③ 行业口径声明（电能量收益的保守下界）
%                ④ 推广的临界条件（价格接受型、不外送、单一储能）
%                ⑤ 结论与建议（不过度声称）
%  禁止空话（如"数据有限""模型有待改进"）；一切改进都必须给出对照数字。
% =============================================================================

\section{模型评价与推广}

\subsection{主要创新点}

每条创新按五段式写：\textbf{缺陷}（残差、约束紧性或对照实验）\(\rightarrow\)\textbf{改动}（可写成公式）
\(\rightarrow\)\textbf{证据}（同一场景集、同一实际轨迹的消融或对照）\(\rightarrow\)\textbf{代价}
\(\rightarrow\)\textbf{边界}（失效条件）。

\begin{innovation}{1}{报童型安全边际：把"该多买多少"写成临界分位数（S 级）}
\ip{缺陷}问题 2 的计划只能依据 \(0{:}00\) 的光伏预报制定，而结算时缺口的代价是正常电价的 \(5\) 倍、
超购只白付电费（两侧成本不对称 \(5:1\)）；按预测中心值制定计划等价于假设预报无误差，但实测逐槽缺货
频率为 \(8.55\%\)、全窗口紧急购电量 \(172\,015.82\)~kWh；确定性阈值律只回答"何时充放电"，
回答不了"计划量该留多少余量"。
\ip{改动}对单槽的超购—缺货权衡建立报童近似：超购成本 \(c_o=p_t\)（超购部分即被弃置），
缺货成本 \(c_u=4p_t\)（紧急购电 \(5p_t\) 中本应付的 \(p_t\) 已计入计划量），故临界分位数
\(F^*(m)=c_u/(c_u+c_o)=4/5\)，即最优缺货概率 \(1/5\)；若超购电量可转存电池并在后续高价时段使用，
取 \(c_o\approx p-p_{\min}=0.7662-0.3713=0.3949\)~元/kWh、\(c_u=4p=3.0648\)~元/kWh 时
\(F^*\approx0.886\)，缺货概率降至约 \(1/10\)。实现上把安全边际写成"按情景缺口分位数超购"的第一阶段量，
第二阶段只对残余缺口紧急购电。
\ip{证据}同一场景集、同一实际轨迹的消融（表~\ref{tab:ablation} 的 E0 行）：移除安全边际
（确定性等价、\(K=1\)）后总费用由 \(14\,264\,792.13\)~元升至 \(17\,398\,559.68\)~元（\(+22.0\%\)），
紧急购电量由 \(172\,015.82\)~kWh 增至 \(1\,156\,618.02\)~kWh（\(6.72\) 倍），且高于无储能基线
\(16\,407\,319.63\)~元（节省率 \(-6.04\%\)）——\textbf{安全边际的边际价值为 \(3\,133\,767.54\)~元}。
方向性检验：\(\kappa\) 由 \(2\) 增到 \(10\) 时安全边际单调升、期望紧急购电单调降（分档数值见
附录表~\ref{tab:sens}），与报童结论方向一致。\evid{C-Q2-01}
\ip{代价}每个情景需解一次第二阶段 LP（单日 \(576+144K\) 个变量，主口径 \(K=60\)），且场景集须先
固定存盘以保证所有变体可比；安全边际使计划购电量偏高、超购部分被弃置（问题 2 窗口盈余处置量
\(3\,106\,390.84\)~kWh）。
\ip{边界}四条失效条件：① 单槽独立近似忽略跨槽储电量耦合与终端条件；② 储能本身即对冲工具，
会替代一部分安全边际（实测缺口频率落在两界下侧）；③ 依赖"紧急购电在结算时刻发生、计划量不可回溯修改"
的时标设定；④ 残差分布未知，只能以样本分位数替代（本文用 \(K=60/120\) 两档量化）。故报童结果只作
\textbf{方向性近似与上下界}，精确解仍由两阶段 LP 给出。
\end{innovation}

\begin{innovation}{2}{预报更新价值的量化：把"要不要加密预报"变成可判定的阈值（S 级）}
\ip{缺陷}题面直接追问"是否需要引入其他时刻的预报"：只答"建议引入更多预报"既不可证伪也无法决策，
只画收益曲线又必然得出"越多越好"。
\ip{改动}在同一固定场景集、同一实际轨迹上\textbf{只改变信息结构}，定义
\(\mathrm{VoI}(n)=\bigl(C(1)-C(n)\bigr)/C(1)\)（\calref{C-10}），并同时给出完全信息上界与"已实现的
可获得信息价值比例"。
\ip{证据}\(C(1)=14\,264\,792.13\)~元 \(\rightarrow\) \(C(2)=14\,131\,564.16\)~元 \(\rightarrow\)
\(C(4)=13\,932\,402.95\)~元，完全信息上界 \(13\,874\,340.39\)~元；\(\mathrm{VoI}(2)=0.93\%\)、
\(\mathrm{VoI}(4)=2.33\%\)，四次预报已实现可获得信息价值的 \(85.13\%\)，再增加预报最多只能再省
\(58\,062.56\)~元（占 \(C(1)\) 的 \(0.42\%\)）；紧急购电量由 \(172\,015.82\)~kWh 降到 \(64\,522.09\)~kWh
（上界 \(6\,134.76\)~kWh）。机理见非预期性置换检验：置换 \(18{:}00\) 预报对全部控制区的影响为 0
（该次预报信息量极低），缺口集中在 \(6{:}00\)--\(12{:}00\)，即若要加密应加在上午而非均匀加密。\evid{C-Q3-01}
\ip{代价}需额外求解两组模型（\(n=2\) 与 \(n=4\)，问题 3 共 \(334\times4=1336\) 个子问题），每天四次
决策带来四次通信、上报与结算复杂度，而全部收益为 \(332\,389.18\)~元（占 \(C(1)\) 的 \(2.33\%\)）。
\ip{边界}结论依赖本题的预报误差结构与惩罚倍数：① 预报误差水平或紧急购电倍数放大时收益平稳带
右移，须重算；② \(18{:}00\) 的预报在本题数据上信息量近于 0，换一套数据后不成立；③ 成本侧未实测——
判据是"单次加密的年化净收益是否为正"，即预报获取的年化成本须低于 \(58\,062.56\)~元；本文未取得气象服务报价，故只给出该阈值而不下"是否值得采购"的结论。
\end{innovation}

\begin{innovation}{3}{波动电价下的联合情景与 CVaR 风险—费用前沿（B 级）}
\ip{缺陷}问题 4 的电价本身随机：实测电价区间 \(0.0076\)--\(1.7936\)~元/kWh、残差组内标准差
\(0.1238\)~元/kWh、一阶自相关 \(0.8589\)，且与净需求正相关（池化 \(0.2602\)）。若对电价、负载、光伏各自
独立抽样，会生成"午间光伏大发同时电价高企"这类不可能情景，从而低估尾部风险；只报期望费用则没有回答"波动"。
\ip{改动}用两层结构建模电价：分时骨架取附件 4 的逐时段均值（与附件 1 的最大偏差
\(5.15\times10^{-5}\)~元/kWh），叠加按日内自相关刻画的随机残差；情景按"日净需求残差 \(\times\) 日均电价"
二维分层取代表历史日（层内 medoid、权重为层样本占比），天然保留电价—净需求相关；风险侧用
Rockafellar--Uryasev 线性化 \(\min\ \mathbb{E}[\text{Cost}]+\lambda\,\mathrm{CVaR}_{0.95}\)
（\calref{C-11}）并对 \(\lambda\) 扫描。
\ip{证据}情景保真度：情景的电价—净需求相关 \(0.2597\) 对历史 \(0.2620\)（差 \(0.0023\)，判据 \(0.05\)），
均值相对差 \(0.082\%\)、关键分位最大相对差 \(6.99\%\)（分块 Bootstrap 对照 \(3.41\%\) 与 \(20.44\%\)）。
前沿（\(56\) 天子样本、日均值）与\textbf{真实轨迹回测}同向：\(\lambda=0\) 时期望 \(41\,995.59\)、
\(\mathrm{CVaR}_{0.95}=52\,176.36\)、回测 \(45\,694.27\) 元/日；\(\lambda=0.5\) 时期望 \(42\,280.55\)
（\(+0.68\%\)）、CVaR \(50\,570.72\)（\(-3.08\%\)）、回测 \(45\,863.46\)（\(+0.37\%\)）、紧急购电量
\(309.62\)~kWh/日（\(-35.52\%\)）；\(\lambda=1\) 时期望 \(42\,411.56\)（\(+0.99\%\)）、CVaR
\(50\,510.46\)（\(-3.19\%\)）、回测 \(45\,952.85\)（\(+0.57\%\)）。即约 \(1\%\) 的期望费用溢价可换取
约 \(3.2\%\) 的尾部风险下降，前沿在 \(\lambda\approx0.6\)--\(0.9\) 后进入平台（\(\lambda=1\) 的 CVaR
反而略高于 \(\lambda=0.9\)）。
\ip{代价}每个 \(\lambda\) 点需重解子样本上的随机模型（共 11 个点），且前沿是"事前"结果、须配真实
轨迹回测才能用于决策；联合情景把电价与负荷抽样耦合成二维分层，情景构造与保真度检验的实现成本上升。
\ip{边界}① 情景集来自 334 天历史，不能外推政策或市场机制突变；
② CVaR 的置信水平固定为 \(0.95\)，未做 \(\alpha\) 扫描；
③ 情景电价截断在实测范围内，出现负电价或极端尖峰时须重新标定；
④ \(\lambda\) 的推荐点取决于决策者风险偏好，本文只给风险—费用前沿，不给唯一答案（\calref{C-11}）。
\end{innovation}

\subsection{模型优点}

\begin{itemize}
  \item \textbf{结构统一}：四问共用同一套物理模型骨架（式~\eqref{eq:q1-balance}--\eqref{eq:q1-nonneg}），
        差别只在信息结构与风险处理，每层新增贡献都可单独消融（表~\ref{tab:ablation} 十组对照）；
  \item \textbf{可解释}：对偶水价值与阈值律给出与求解器无关的策略形式（\(\alpha_t\le\eta_c\nu_t\) 充电、
        \(\alpha_t\ge\nu_t/\eta_d\) 放电，全部 144 个槽违反次数 0），使充放电时段可解释、可质疑；
  \item \textbf{可复算}：全部数值由脚本从原始附件重算，并通过原始可行性（供给缺口 \(3.41\times10^{-13}\)~kWh）、
        最优性（阈值律一致性）与费用独立复算（相对误差 0）三重检验（第九章）。
\end{itemize}

\subsection{模型局限}



\textbf{本文目标函数只含电能量购电费用，未计入容量补偿、辅助服务收益与绿电环境价值，
也未计入电池循环衰减与安全成本；因此本文所得节约幅度应视为储能电能量价值的保守下界。}
用户到户电价还含输配电费、系统运行费与政府性基金及附加，本文只优化第一项：附加电费按电量比例收取
等价于对全部电价加同一常数，最优调度会因往返效率小于 1 而改变；按需量（容量）收取则更会实质性改变
策略。两种情形都须重解，故本文的节省幅度既不是储能项目投资可行性的结论，也不是全成本口径的收益。
按电量比例计收的附加费等价于给全部电价加上同一常数，由于往返效率小于 1，最优充放电量会相应减少（本文未重解，故不给出该口径下的数值）；按最大需量或容量计收的附加费则会实质性改变策略，须在目标函数中显式建模后重解。

\subsection{模型推广与临界条件}

本文模型在以下条件下成立，偏离时须重新校准而非直接套用：

\begin{itemize}
  \item \textbf{价格接受型}：电价外生给定（附件 1 与附件 4）。若储能或微网规模足以影响出清价格，
        目标函数须引入价格对自身报价的响应，模型由 LP 变为带均衡约束的问题，本文的阈值律与安全边际
        结论都不能直接沿用；
  \item \textbf{不外送（假设 A1）}：允许按同一电价售电时，目标函数需增加售电收益项，盈余处置量
        \calref{C-06} 由"必要"变为"可选"，收益结构整体改变（售电可能成为主要收益来源），须重解；
  \item \textbf{单一储能}：推广到多储能或多微网协同后状态维度随设备数增长，"逐日解耦"的代价须
        重新评估——本文数据上为 \(1\,183\,394.34\)~元（\(9.68\%\)，表~\ref{tab:ablation} 的 FY 行）；
  \item \textbf{价差水平足够大}：套利有利的盈亏平衡条件是谷峰价比高于 \(1/(\eta_c\eta_d)=1.2346\)
        （主口径单边各 \(0.9\)）；价差压到该阈值附近时纯套利不再有利，模型退化为"以规避紧急购电为主"
        的保险型策略。容量与功率扫描给出同类临界：容量降到 \(7\,200\)~kWh 时节省率由 \(26.90\%\) 降到
        \(19.59\%\)，功率降到 \(1\,250\)~kW 时降到 \(16.41\%\)；
  \item \textbf{负载刚性}：若引入需求响应，式~\eqref{eq:q1-balance} 中的 \(L_t\) 变为决策量，可行域
        变为含柔性负荷约束的多面体，阈值律须重新推导（水价值与负荷转移的边际价值共同决定阈值），
        但"两阶段 + 安全边际"框架可沿用；
  \item \textbf{单一电能量结算}：容量补偿、辅助服务与绿电环境价值在多数地区顺序出清、难以同时兑现，
        收益\textbf{不可简单叠加}；计入后本文结论方向不变（当前口径为保守下界），但幅度须重算。
\end{itemize}

\textbf{结论与建议。}在本题参数范围内（容量 \(12\,000\)~kWh、充放功率 \(5000\)~kW、单边效率 \(0.9\)、
紧急购电 \(5\) 倍价、不外送），储能的价值是把低谷电量与午间光伏盈余搬运到高价时段：典型日节省
\(26.90\%\)，334 天窗口在预报不确定下降到 \(13.06\%\)，波动电价下为 \(11.35\%\)；容量摆幅是主要瓶颈
（仅解除容量上限可再省 \(22.14\%\)，仅解除功率上限再省 \(1.66\%\)）；预报价值已接近饱和（四次预报实现
可获得信息价值的 \(85.13\%\)，再加密的收益上限为 \(58\,062.56\)~元）；风险偏好可通过 \(\lambda\) 调节，
约 \(1\%\) 的期望费用溢价可换约 \(3.2\%\) 的尾部风险下降。据此给出三条建议：① 预算有限时优先扩容储能
容量而不是加密预报（容量档位的收益区间为 \(19.59\%\)--\(42.56\%\)，远大于加密预报的上限 \(0.42\%\)）；
② 安全边际必须保留（边际价值 \(3\,133\,767.54\)~元），并应加在上午的缺口高发时段；③ 部署前按
表~\ref{tab:limitation} 逐项校准效率、衰减、负载预测与终端条件。上述结论只对"价格接受型、不外送、
单一储能、以电能量购电费用最小"这一设定成立，不构成储能项目的投资可行性结论。
